\paragraph{Computation of Sample-Level Blink-Region Threshold Candidates}

Stage~B estimates channel-specific sample-level amplitude thresholds using signal samples drawn from the candidate epochs \(\mathcal{B}\) identified in Stage~A.
The estimation combines a measure of central tendency with a robust measure of dispersion and a sensitivity parameter, allowing the threshold to adapt to the amplitude characteristics of each channel while limiting the influence of isolated extreme samples.

%In Stage~B, a robust sample-level amplitude threshold is estimated using only the samples
%pooled from the flagged epochs \[\mathcal{B}\], as discussed in previous subsection., as
%\(\theta = \xi + k\,(1.4826 \cdot \mathrm{MAD})\), where the centre \(\xi\) is the
%\emph{median} in the primary configuration (or the \emph{mean} in the contrast
%configuration), \(\mathrm{MAD}\) is the median absolute deviation, the factor
%\(1.4826\) rescales the MAD to a standard-deviation-consistent dispersion, and \(k\) is
%the sensitivity multiplier. Estimating the threshold from the flagged epochs concentrates
%the calculation on blink-enriched signal, while the median-plus-MAD construction keeps the
%estimate stable in the presence of isolated extreme samples. If too few epochs are flagged,
%the estimate falls back to the full valid-epoch set so that the threshold is never derived
%from an unreliably small sample.
%
%
%The second stage converts the epoch-level artifact decision into a
%sample-level thresholding problem. The autoreject-style screening in the
%previous stage identifies a set of suspicious epochs \(\mathcal{B}\), but its
%thresholds \(\tau_c^\star\) are defined on peak-to-peak amplitudes over whole
%epochs. Such thresholds are therefore appropriate for deciding whether an epoch
%is abnormal, but they are not directly interpretable as sample-wise blink
%crossing levels. Stage B estimates a separate channel-specific threshold in
%the amplitude domain of the time series itself.

Let \(\mathcal{S}\) denote the set of epochs used for estimating the
sample-level threshold. In the primary analysis,
\(\mathcal{S}=\mathcal{B}\), so that the estimate is driven by epochs already
identified as likely to contain blink-related activity.

%If too few suspicious
%epochs are available, the method falls back to the full valid epoch set
%\(\mathcal{E}\) to avoid estimating a threshold from an insufficient sample:
%\[
%\mathcal{S}
%=
%\begin{cases}
%\mathcal{B}, & |\mathcal{B}| \geq n_{\min},\\
%\mathcal{E}, & |\mathcal{B}| < n_{\min},
%\end{cases}
%\]
%where \(n_{\min}\) is the minimum number of screened epochs required for a
%stable estimate.

For each channel \(c\), all samples from the selected epochs are pooled into a
one-dimensional empirical sample set
\[
\mathcal{Y}_c
= \{x_{i,c,t}: i \in \mathcal{S},\; 1 \leq t \leq T\}.
\]
%This pooling step marks the methodological transition from epoch-level rejection to sample-level localization: the quantities used in Stage B are no
%longer peak-to-peak ranges \(a_{i,c}\), but the individual signal amplitudes
%\(x_{i,c,t}\) that will later be thresholded to form candidate blink regions.

The sample-level threshold is estimated using a robust location-plus-dispersion form. Let \(\xi_c\) be the channel-specific center of the sample distribution.
The default choice is the median,
\[
\xi_c = \operatorname{median}_{y \in \mathcal{Y}_c} y,
\]

It should be note that, the arithmetic mean may be substituted when a mean-centered threshold
is required for comparison.

The dispersion is estimated by the median absolute
deviation (MAD),
\[
\operatorname{MAD}_c
= \operatorname{median}_{y \in \mathcal{Y}_c}
\left|y-\operatorname{median}_{u \in \mathcal{Y}_c}u\right|,
\]
and is rescaled to a standard-deviation-compatible quantity by
\[
\sigma^{\mathrm{rob}}_c = 1.4826\,\operatorname{MAD}_c .
\]
The constant \(1.4826\) is the usual normal-consistency factor for the MAD.

For a sensitivity parameter \(k>0\), the resulting sample-level blink-region
threshold candidate for channel \(c\) is
\[
\theta_c = \xi_c + k\sigma^{\mathrm{rob}}_c .
\]
The set
\[
\Theta = \{\theta_c: c \in \mathcal{C}\}
\]
constitutes the channel-wise threshold candidates passed to the blink-region extraction stage.

%This construction deliberately separates the statistical role of the two thresholds in the pipeline. The quantities \(\tau_c^\star\) from Stage A are  epoch-rejection thresholds defined on peak-to-peak amplitudes,
%whereas
%\(\theta_c\) is a sample-level amplitude threshold defined on the time-series  samples themselves.
%Estimating \(\theta_c\) from \(\mathcal{S}\) concentrates  the calculation on signal segments enriched for blink-like excursions, while  the robust center and MAD-scaled dispersion reduce the influence of isolated  extreme peaks within those same segments.
